\section{Spray Physics \& Atomization}

\subsection{Introduction}
This chapter describes the physics of propellant atomization, spray formation, and evaporation. The models track the transformation of liquid propellant streams into droplets and their subsequent gasification.

\subsection{Key Files Referenced}
The physics models documented in this chapter are implemented in:
\begin{itemize}
    \item \texttt{engine/core/spray.py} - Spray and atomization models
    \item \texttt{engine/core/injectors/pintle.py} - Spray integration for pintle injectors
\end{itemize}

\subsection{Nomenclature}
\begin{description}
    \item[$J$] Momentum flux ratio [dimensionless]
    \item[$TMR$] Thrust momentum ratio [dimensionless]
    \item[$\theta$] Spray angle [\si{rad}]
    \item[$We$] Weber number [dimensionless]
    \item[$Oh$] Ohnesorge number [dimensionless]
    \item[$D_{32}$] Sauter Mean Diameter (SMD) [\si{m}]
    \item[$\tau_{evap}$] Evaporation time [\si{s}]
    \item[$x^*$] Evaporation length [\si{m}]
    \item[$\sigma$] Surface tension [\si{N/m}]
\end{description}

\subsection{Momentum and Spray Angle}

\subsubsection{Momentum Flux Ratio}
The momentum flux ratio $J$ characterizes the relative strength of the oxidizer and fuel streams:
\begin{equation}
J = \frac{\rho_{ox} u_{ox}^2}{\rho_{fuel} u_{fuel}^2}
\end{equation}

\subsubsection{Thrust Momentum Ratio (TMR)}
For pintle injectors, the TMR is defined as:
\begin{equation}
TMR = \frac{\mdot_{fuel} u_{fuel}}{\mdot_{ox} u_{ox}}
\end{equation}

\subsubsection{Spray Angle}
The spray angle $\theta$ is modeled as a function of the momentum ratio:
\begin{equation}
\theta = \arctan(a_J \sqrt{J})
\end{equation}
Alternatively, using TMR:
\begin{equation}
\theta = \arccos\left(\frac{1}{1 + TMR}\right)
\end{equation}

\subsection{Atomization (SMD Models)}

\subsubsection{Weber Number}
The Weber number characterizes droplet breakup:
\begin{equation}
We = \frac{\rho_{gas} u_{rel}^2 D}{\sigma}
\end{equation}

\subsubsection{Ohnesorge Number}
The Ohnesorge number accounts for viscous effects:
\begin{equation}
Oh = \frac{\mu_{liq}}{\sqrt{\rho_{liq} \sigma D}}
\end{equation}

\subsubsection{Sauter Mean Diameter (SMD)}
The SMD is modeled using the Lefebvre correlation for pressure-atomized sprays:
\begin{equation}
D_{32} = a_1 \sigma^{0.5} \mu_{liq}^{0.1} \mdot^{0.1} \Delta P^{-0.5} \rho_{gas}^{-0.25}
\end{equation}
For pintle-specific geometries:
\begin{equation}
D_{32, pintle} = a_2 \left(\frac{\sigma}{\rho_{gas} u_{rel}^2}\right)^{0.5} \left(\frac{h_{gap}}{d_{orifice}}\right)^{0.25}
\end{equation}

\subsection{Evaporation and Gasification}

\subsubsection{Evaporation Time}
The characteristic time for droplet evaporation is:
\begin{equation}
\tau_{evap} = \frac{D_{32}^2}{K_{evap}}
\end{equation}
where $K_{evap}$ is the evaporation constant.

\subsubsection{Evaporation Length}
The axial distance required for complete evaporation:
\begin{equation}
x^* = u_{drop} \tau_{evap}
\end{equation}

\subsection{Pintle-Specific Spray Integration}
For the pintle injector, the spray angle and atomization quality are coupled to the annular gap height and the radial momentum of the fuel sheet.
